% !PRINTED 536
\weberpage{536}
%<<CONT
Ruhe lässt. Ist dann $\pi_i$ eine Permutation, die die Ziffer 0 in $i$ überführt, so ist $\pi_i^{-1}\,Q\,\pi_i$ die Gruppe der Permutationen in $P$, die die Ziffer $i$ nicht ändern. Da nun aber in jeder Zusammensetzung von Permutationen aus $P$ die Componenten vertauscht werden können, so ist
\[
\pi_i^{-1}\,Q\,\pi_i = \pi_i^{-1}\pi_i\,Q = Q,
\]
d.\,h. die Gruppe $Q$ lässt auch die Ziffer $i$ ungeändert. Wegen der Transitivität von $P$ kann aber $i$ jede der Ziffern $1,2\ldots m-1$ bedeuten, folglich besteht $Q$ aus der einzigen identischen Permutation, und es giebt in $P$ ausser der identischen keine Permutation, die eine Ziffer ungeändert lässt. Adjungirt man aber \gesperrt{eine} Wurzel $\alpha$, so reducirt sich die Gruppe $P$ auf die Einheitsgruppe oder die Gleichung ist gelöst; $\varphi(x)=0$ ist eine Normalgleichung und somit ihre eigene \gesperrt{Galois'sche} Resolvente. Ist $\alpha_k=\Theta_k(\alpha)$, so besteht die \gesperrt{Galois'sche} Gruppe dieser Gleichung aus den Substitutionen:
\[
\sigma_k = [\alpha,\Theta_k(\alpha)],
\]
und es ist
\[
% [sic] authority prints x in the second composite; the parallel expression and the next line require \alpha.
\sigma_h\sigma_k = [\alpha,\Theta_h\Theta_k(\alpha)],\quad \sigma_k\sigma_h = [\alpha,\Theta_k\Theta_h(x)].
\]
Da die Gruppe commutativ sein soll, so muss
\[
\Theta_h\Theta_k(\alpha) = \Theta_k\Theta_h(\alpha)
\]
sein, d.\,h. $\varphi(x)=0$ ist eine \gesperrt{Abel'sche Gleichung}. Dies ist der Grund, weshalb die commutativen Gruppen auch \gesperrt{Abel'sche} Gruppen genannt werden (\S.\,148). Es folgt also noch aus dem oben gegebenen Satze über die Gruppe von Normalgleichungen, dass bei einer transitiven \gesperrt{Abel'schen} Permutationsgruppe die Zahl der Permutationen mit der Zahl der vertauschten Ziffern übereinstimmt.

Hier haben wir die Irreducibilität der gegebenen Gleichung vorausgesetzt. Wollen wir auch reducible Gleichungen mit commutativer Gruppe in Betracht ziehen, so brauchen wir nur an Stelle der Gleichung selbst eine ihrer \gesperrt{Galois'schen} Resolventen zu setzen.

Es ist endlich noch zu bemerken, dass bei einer commutativen Gruppe jeder Theiler normal ist, da ja immer $\pi^{-1}\varkappa\pi = \varkappa$ ist, wenn $\pi$ und $\varkappa$ beliebige Elemente einer commutativen Gruppe sind.
